\documentclass[10pt,a4paper]{amsart}
\usepackage{amsmath,amssymb,amsthm,graphicx}
% Retained preamble for the Stacks Project chapter snapshot in this directory.
% The chapter's own preamble.tex is a copy of the shared Stacks Project preamble; the block
% below is sliced MECHANICALLY, line for line, from it. Nothing was added, removed or
% redefined; the wrapper class, the xr-hyper/hyperref wiring and the externaldocument list
% are omitted because the assembler supplies the document class.
%
% lines taken from preamble.tex: 3, 5, 9, 11, 12, 156, 157, 158, ... (51 lines total)
%
% package coverage restored from the chapter's own preamble (2026-10-13): xy, xy-2cell

\usepackage[all]{xy}
\xyoption{2cell}
\UseAllTwocells

\newenvironment{reference}{\comment}{\endcomment}
\newenvironment{slogan}{\comment}{\endcomment}
\newenvironment{history}{\comment}{\endcomment}
\theoremstyle{plain}
\newtheorem{theorem}[subsection]{Theorem}
\newtheorem{proposition}[subsection]{Proposition}
\newtheorem{lemma}[subsection]{Lemma}
\theoremstyle{definition}
\newtheorem{definition}[subsection]{Definition}
\newtheorem{example}[subsection]{Example}
\newtheorem{exercise}[subsection]{Exercise}
\newtheorem{situation}[subsection]{Situation}
\theoremstyle{remark}
\newtheorem{remark}[subsection]{Remark}
\newtheorem{remarks}[subsection]{Remarks}
\numberwithin{equation}{subsection}
\def\lim{\mathop{\mathrm{lim}}\nolimits}
\def\colim{\mathop{\mathrm{colim}}\nolimits}
\def\Spec{\mathop{\mathrm{Spec}}}
\def\Hom{\mathop{\mathrm{Hom}}\nolimits}
\def\Ext{\mathop{\mathrm{Ext}}\nolimits}
\def\SheafHom{\mathop{\mathcal{H}\!\mathit{om}}\nolimits}
\def\SheafExt{\mathop{\mathcal{E}\!\mathit{xt}}\nolimits}
\def\Sch{\mathit{Sch}}
\def\Mor{\mathop{\mathrm{Mor}}\nolimits}
\def\Ob{\mathop{\mathrm{Ob}}\nolimits}
\def\Sh{\mathop{\mathit{Sh}}\nolimits}
\def\NL{\mathop{N\!L}\nolimits}
\def\CH{\mathop{\mathrm{CH}}\nolimits}
\def\proetale{{pro\text{-}\acute{e}tale}}
\def\etale{{\acute{e}tale}}
\def\QCoh{\mathit{QCoh}}
\def\Ker{\mathop{\mathrm{Ker}}}
\def\Im{\mathop{\mathrm{Im}}}
\def\Coker{\mathop{\mathrm{Coker}}}
\def\Coim{\mathop{\mathrm{Coim}}}
\DeclareMathSymbol{\boxtimes}{\mathbin}{AMSa}{"02}
\def\QCohstack{\mathcal{QC}\!\mathit{oh}}
\def\Cohstack{\mathcal{C}\!\mathit{oh}}
\def\Spacesstack{\mathcal{S}\!\mathit{paces}}
\def\Quotfunctor{\mathrm{Quot}}
\def\Hilbfunctor{\mathrm{Hilb}}
\def\Curvesstack{\mathcal{C}\!\mathit{urves}}
\def\Polarizedstack{\mathcal{P}\!\mathit{olarized}}
\def\Complexesstack{\mathcal{C}\!\mathit{omplexes}}
\def\Pic{\mathop{\mathrm{Pic}}\nolimits}
\def\Picardstack{\mathcal{P}\!\mathit{ic}}
\def\Picardfunctor{\mathrm{Pic}}
\def\Deformationcategory{\mathcal{D}\!\mathit{ef}}
\renewenvironment{reference}{}{}
\renewenvironment{slogan}{}{}
\renewenvironment{history}{}{}
\usepackage{cleveref}
\crefname{originalproposition}{original proposition}{original propositions}
\Crefname{originalproposition}{Original proposition}{Original propositions}
\makeatletter
\newlabel{algebra-lemma-support-dimension-d}{{00L8}{source}{}{sourcepointer-1}{}}
\newlabel{algebra-lemma-support-dimension-d@cref}{{[originalproposition][0][]00L8}{[1][0][]source}{}{}{}}
\makeatother
\crefname{originalproposition}{original proposition}{original propositions}
\Crefname{originalproposition}{Original proposition}{Original propositions}
\makeatletter
\newlabel{dualizing-lemma-special-case-local-duality}{{0AAK}{source}{}{sourcepointer-2}{}}
\newlabel{dualizing-lemma-special-case-local-duality@cref}{{[originalproposition][0][]0AAK}{[1][0][]source}{}{}{}}
\makeatother
\crefname{originalproposition}{original proposition}{original propositions}
\Crefname{originalproposition}{Original proposition}{Original propositions}
\makeatletter
\newlabel{dualizing-lemma-universally-catenary}{{0A80}{source}{}{sourcepointer-3}{}}
\newlabel{dualizing-lemma-universally-catenary@cref}{{[originalproposition][0][]0A80}{[1][0][]source}{}{}{}}
\makeatother
\crefname{originalproposition}{original proposition}{original propositions}
\Crefname{originalproposition}{Original proposition}{Original propositions}
\makeatletter
\newlabel{lemma-nonnegative-chi-triple}{{0F28}{source}{}{sourcepointer-4}{}}
\newlabel{lemma-nonnegative-chi-triple@cref}{{[originalproposition][0][]0F28}{[1][0][]source}{}{}{}}
\makeatother
\crefname{originalproposition}{original proposition}{original propositions}
\Crefname{originalproposition}{Original proposition}{Original propositions}
\makeatletter
\newlabel{local-cohomology-lemma-catenary-S2-equidimensional}{{0FIW}{source}{}{sourcepointer-5}{}}
\newlabel{local-cohomology-lemma-catenary-S2-equidimensional@cref}{{[originalproposition][0][]0FIW}{[1][0][]source}{}{}{}}
\makeatother
\crefname{originalproposition}{original proposition}{original propositions}
\Crefname{originalproposition}{Original proposition}{Original propositions}
\makeatletter
\newlabel{local-cohomology-lemma-local-finiteness-for-finite-locally-free}{{0BPY}{source}{}{sourcepointer-6}{}}
\newlabel{local-cohomology-lemma-local-finiteness-for-finite-locally-free@cref}{{[originalproposition][0][]0BPY}{[1][0][]source}{}{}{}}
\makeatother
\crefname{originalproposition}{original proposition}{original propositions}
\Crefname{originalproposition}{Original proposition}{Original propositions}
\makeatletter
\newlabel{local-cohomology-lemma-sitting-in-degrees}{{0DWZ}{source}{}{sourcepointer-7}{}}
\newlabel{local-cohomology-lemma-sitting-in-degrees@cref}{{[originalproposition][0][]0DWZ}{[1][0][]source}{}{}{}}
\makeatother
\crefname{originalproposition}{original proposition}{original propositions}
\Crefname{originalproposition}{Original proposition}{Original propositions}
\makeatletter
\newlabel{properties-lemma-quasi-affine-invertible-nonvanishing-section}{{0F20}{source}{}{sourcepointer-8}{}}
\newlabel{properties-lemma-quasi-affine-invertible-nonvanishing-section@cref}{{[originalproposition][0][]0F20}{[1][0][]source}{}{}{}}
\makeatother
\crefname{originalproposition}{original proposition}{original propositions}
\Crefname{originalproposition}{Original proposition}{Original propositions}
\makeatletter
\newlabel{proposition-hilbert-triple}{{0F27}{source}{}{sourcepointer-9}{}}
\newlabel{proposition-hilbert-triple@cref}{{[originalproposition][0][]0F27}{[1][0][]source}{}{}{}}
\makeatother
\crefname{originalproposition}{original proposition}{original propositions}
\Crefname{originalproposition}{Original proposition}{Original propositions}
\makeatletter
\newlabel{section-coherent-triples}{{0F22}{source}{}{sourcepointer-10}{}}
\newlabel{section-coherent-triples@cref}{{[originalproposition][0][]0F22}{[1][0][]source}{}{}{}}
\makeatother
\title{Kollar: under depth and dimension hypotheses the Picard map from the punctured spectrum to the punctured spectrum mod f is injective}
\author{The Stacks Project Authors}
\date{stacks-project:algebraization.tex:local-snapshot-20261012; GNU FDL 1.2 or later}
\begin{document}
\maketitle
\noindent\textit{Editorial note.} Statement span is lines 8729-8749 of sources/stacks-algebraization/pinned-original-source.tex; proof span is lines 8751-8904, the complete author proof (152 body lines).\par\medskip
\noindent\textit{Editorial note.} Navigation stubs: 10 cited-but-not-included labels are supplied as non-load-bearing pointers with their pinned tags where the snapshot tag index has them.\par\medskip
\noindent\textit{Editorial note.} extra\_packages is empty on purpose: the retained preamble already loads xy with options and a second declaration clashes in the assembler wrapper.\par\medskip
\noindent\textit{Editorial note.} Text is retained verbatim, including the author's own typos and placeholders.\par\medskip
\noindent\textit{Editorial note.} Pool algebraization wave36 (al02 bib-salvage batch). Spans computed from \textbackslash{}label (16.5/16.6).\par\medskip
\noindent\textit{Editorial note.} External \textbackslash{}cite keys salvaged with card-local bibliography.tex entries copied verbatim from the snapshot my.bib (precedent S3-morph-blow-up-ample-family); URL/ISSN/EPRINT fields dropped as the preamble has no url package.\par\medskip
\noindent\textit{Editorial note.} statement contains reference environment with real content -> patch-sd-envs.tex applied (sd02 lesson).\par\medskip
\noindent\textit{Editorial note.} Cross-chapter same-label warning (NOT an occupancy conflict, Chief ruling 2026-10-13 evening): label 'proposition-injective-pic' also exists in stacks-snapshot-20261011 — a different proposition in a different chapter file.\par\medskip
\noindent\textit{Allowed standard prerequisites (editorial).}\begin{itemize}
\item Explicit prerequisite boundary. See why-field; in-chapter lemmas cited by name become stub prereqs; external cites resolved via the card-local bibliography.tex fragment.
\end{itemize}\medskip
\ifdefined\hypertarget\hypertarget{sourcepointer-1}{}\fi
\noindent\textit{Source reference.} Original proposition 00L8: sources/stacks-algebraization/pinned-original-source.tex, source lines 8904--8904. Cited by the retained statement or proof but not part of the collected spans; supplied only as a navigation pointer. The argument does not depend on its content.\par\medskip
\ifdefined\hypertarget\hypertarget{sourcepointer-2}{}\fi
\noindent\textit{Source reference.} Original proposition 0AAK: sources/stacks-algebraization/pinned-original-source.tex, source lines 8904--8904. Cited by the retained statement or proof but not part of the collected spans; supplied only as a navigation pointer. The argument does not depend on its content.\par\medskip
\ifdefined\hypertarget\hypertarget{sourcepointer-3}{}\fi
\noindent\textit{Source reference.} Original proposition 0A80: sources/stacks-algebraization/pinned-original-source.tex, source lines 8904--8904. Cited by the retained statement or proof but not part of the collected spans; supplied only as a navigation pointer. The argument does not depend on its content.\par\medskip
\ifdefined\hypertarget\hypertarget{sourcepointer-4}{}\fi
\noindent\textit{Source reference.} Original proposition 0F28: sources/stacks-algebraization/pinned-original-source.tex, source lines 8612--8623. Cited by the retained statement or proof but not part of the collected spans; supplied only as a navigation pointer. The argument does not depend on its content.\par\medskip
\ifdefined\hypertarget\hypertarget{sourcepointer-5}{}\fi
\noindent\textit{Source reference.} Original proposition 0FIW: sources/stacks-algebraization/pinned-original-source.tex, source lines 8904--8904. Cited by the retained statement or proof but not part of the collected spans; supplied only as a navigation pointer. The argument does not depend on its content.\par\medskip
\ifdefined\hypertarget\hypertarget{sourcepointer-6}{}\fi
\noindent\textit{Source reference.} Original proposition 0BPY: sources/stacks-algebraization/pinned-original-source.tex, source lines 8904--8904. Cited by the retained statement or proof but not part of the collected spans; supplied only as a navigation pointer. The argument does not depend on its content.\par\medskip
\ifdefined\hypertarget\hypertarget{sourcepointer-7}{}\fi
\noindent\textit{Source reference.} Original proposition 0DWZ: sources/stacks-algebraization/pinned-original-source.tex, source lines 8904--8904. Cited by the retained statement or proof but not part of the collected spans; supplied only as a navigation pointer. The argument does not depend on its content.\par\medskip
\ifdefined\hypertarget\hypertarget{sourcepointer-8}{}\fi
\noindent\textit{Source reference.} Original proposition 0F20: sources/stacks-algebraization/pinned-original-source.tex, source lines 8904--8904. Cited by the retained statement or proof but not part of the collected spans; supplied only as a navigation pointer. The argument does not depend on its content.\par\medskip
\ifdefined\hypertarget\hypertarget{sourcepointer-9}{}\fi
\noindent\textit{Source reference.} Original proposition 0F27: sources/stacks-algebraization/pinned-original-source.tex, source lines 8503--8514. Cited by the retained statement or proof but not part of the collected spans; supplied only as a navigation pointer. The argument does not depend on its content.\par\medskip
\ifdefined\hypertarget\hypertarget{sourcepointer-10}{}\fi
\noindent\textit{Source reference.} Original proposition 0F22: sources/stacks-algebraization/pinned-original-source.tex, source lines 8168--8183. Cited by the retained statement or proof but not part of the collected spans; supplied only as a navigation pointer. The argument does not depend on its content.\par\medskip
\begin{proposition}[Koll\'ar]
\label{proposition-injective-pic}
\begin{reference}
\cite[Theorem 1.9]{Kollar-map-pic}
\end{reference}
Let $(A, \mathfrak m)$ be a Noetherian local ring. Let $f \in \mathfrak m$.
Assume
\begin{enumerate}
\item $A$ has a dualizing complex,
\item $f$ is a nonzerodivisor,
\item $\text{depth}(A/fA) \geq 2$, or equivalently $\text{depth}(A) \geq 3$,
\item if $f \in \mathfrak p \subset A$ is a prime ideal with
$\dim(A/\mathfrak p) = 2$, then $\text{depth}(A_\mathfrak p) \geq 2$.
\end{enumerate}
Let $U$, resp.\ $U_0$ be the punctured spectrum of $A$, resp.\ $A/fA$. The map
$$
\Pic(U) \to \Pic(U_0)
$$
is injective. Finally, if (1), (2), (3), $A$ is $(S_2)$, and
$\dim(A) \geq 4$, then (4) holds.
\end{proposition}

\begin{proof}
Let $\mathcal{L}$ be an invertible $\mathcal{O}_U$-module.
Observe that $\mathcal{L}$ maps to $0$ in $\Pic(U_0)$
if and only if we can extend $\mathcal{L}$ to an invertible
coherent triple $(\mathcal{L}, \mathcal{L}_0, \lambda)$
as in Section \ref{section-coherent-triples}.
By Proposition \ref{proposition-hilbert-triple}
the function
$$
n \longmapsto \chi((\mathcal{L}, \mathcal{L}_0, \lambda)^{\otimes n})
$$
is a polynomial $P$. By Lemma \ref{lemma-nonnegative-chi-triple}
we have
$P(n) \geq 0$ for all $n \in \mathbf{Z}$ with equality if and only if
$\mathcal{L}^{\otimes n}$ is trivial. In particular $P(0) = 0$
and $P$ is either identically zero and we win or $P$ has even degree $\geq 2$.

\medskip\noindent
Set $M = \Gamma(U, \mathcal{L})$ and
$M_0 = \Gamma(X_0, \mathcal{L}_0) = \Gamma(U_0, \mathcal{L}_0)$.
Then $M$ is a finite $A$-module of depth $\geq 2$
and $M_0 \cong A/fA$, see proof of Lemma \ref{lemma-nonnegative-chi-triple}.
Note that $H^2_\mathfrak m(M)$ is finite $A$-module by
Local Cohomology, Lemma
\ref{local-cohomology-lemma-local-finiteness-for-finite-locally-free}
and the fact that $H^i_\mathfrak m(A) = 0$ for $i = 0, 1, 2$
since $\text{depth}(A) \geq 3$.
Consider the short exact sequence
$$
0 \to M/fM \to M_0 \to Q \to 0
$$
Lemma \ref{lemma-nonnegative-chi-triple} tells us $Q$ has finite length
equal to $\chi(\mathcal{L}, \mathcal{L}_0, \lambda)$.
We obtain $Q = H^1_\mathfrak m(M/fM)$ and
$H^i_\mathfrak m(M/fM) = H^i_\mathfrak m(M_0) \cong H^i_\mathfrak m(A/fA)$
for $i > 1$ from the long exact sequence of local cohomology
associated to the displayed short exact sequence. Consider the long
exact sequence of local cohomology associated to the sequence
$0 \to M \to M \to M/fM \to 0$. It starts with
$$
0 \to Q \to H^2_\mathfrak m(M) \to H^2_\mathfrak m(M) \to
H^2_\mathfrak m(A/fA)
$$
Using additivity of lengths we see that
$\chi(\mathcal{L}, \mathcal{L}_0, \lambda)$
is equal to the length of the image of
$H^2_\mathfrak m(M) \to H^2_\mathfrak m(A/fA)$.

\medskip\noindent
Let prove the lemma in a special case to elucidate the rest of the proof.
Namely, assume for a moment that $H^2_\mathfrak m(A/fA)$ is
a finite length module. Then
we would have $P(1) \leq \text{length}_A H^2_\mathfrak m(A/fA)$.
The exact same argument applied to $\mathcal{L}^{\otimes n}$ shows that
$P(n) \leq \text{length}_A H^2_\mathfrak m(A/fA)$ for all $n$.
Thus $P$ cannot have positive degree and we win.
In the rest of the proof we will modify this argument to give
a linear upper bound for $P(n)$ which suffices.

\medskip\noindent
Let us study the map
$H^2_\mathfrak m(M) \to H^2_\mathfrak m(M_0) \cong H^2_\mathfrak m(A/fA)$.
Choose a normalized dualizing complex $\omega_A^\bullet$ for $A$.
By local duality
(Dualizing Complexes, Lemma \ref{dualizing-lemma-special-case-local-duality})
this map is Matlis dual to the map
$$
\text{Ext}^{-2}_A(M, \omega_A^\bullet)
\longleftarrow
\text{Ext}^{-2}_A(M_0, \omega_A^\bullet)
$$
whose image therefore has the same (finite) length.
The support (if nonempty) of the finite $A$-module
$\text{Ext}^{-2}_A(M_0, \omega_A^\bullet)$ consists of
$\mathfrak m$ and a finite number of primes
$\mathfrak p_1, \ldots, \mathfrak p_r$ containing $f$ with
$\dim(A/\mathfrak p_i) = 1$. Namely, by
Local Cohomology, Lemma \ref{local-cohomology-lemma-sitting-in-degrees}
the support is contained in the set of primes $\mathfrak p \subset A$ with
$\text{depth}_{A_\mathfrak p}(M_{0, \mathfrak p}) + \dim(A/\mathfrak p) \leq 2$.
Thus it suffices to show there is no prime $\mathfrak p$ containing $f$ with
$\dim(A/\mathfrak p) = 2$ and
$\text{depth}_{A_\mathfrak p}(M_{0, \mathfrak p}) = 0$.
However, because $M_{0, \mathfrak p} \cong (A/fA)_\mathfrak p$
this would give $\text{depth}(A_\mathfrak p) = 1$ which contradicts
assumption (4).
Choose a section $t \in \Gamma(U, \mathcal{L}^{\otimes -1})$
which does not vanish in the points $\mathfrak p_1, \ldots, \mathfrak p_r$, see
Properties, Lemma
\ref{properties-lemma-quasi-affine-invertible-nonvanishing-section}.
Multiplication by $t$ on global sections determines a map $t : M \to A$
which defines an isomorphism $M_{\mathfrak p_i} \to A_{\mathfrak p_i}$ for
$i = 1, \ldots, r$. Denote $t_0 = t|_{U_0}$ the corresponding section
of $\Gamma(U_0, \mathcal{L}_0^{\otimes -1})$ which similarly determines
a map $t_0 : M_0 \to A/fA$ compatible with $t$.
We conclude that there is a commutative diagram
$$
\xymatrix{
\text{Ext}^{-2}_A(M, \omega_A^\bullet) &
\text{Ext}^{-2}_A(M_0, \omega_A^\bullet) \ar[l] \\
\text{Ext}^{-2}_A(A, \omega_A^\bullet) \ar[u]^t &
\text{Ext}^{-2}_A(A/fA, \omega_A^\bullet) \ar[l] \ar[u]_{t_0}
}
$$
It follows that the length of the image of the top horizontal
map is at most the length of $\text{Ext}^{-2}_A(A/fA, \omega_A^\bullet)$
plus the length of the cokernel of $t_0$.

\medskip\noindent
However, if we replace $\mathcal{L}$ by $\mathcal{L}^n$ for $n > 1$,
then we can use
$$
t^n :
M_n = \Gamma(U, \mathcal{L}^{\otimes n})
\longrightarrow
\Gamma(U, \mathcal{O}_U) = A
$$
instead of $t$. This replaces $t_0$ by its $n$th power.
Thus the length of the image of the map
$\text{Ext}^{-2}_A(M_n, \omega_A^\bullet) \leftarrow
\text{Ext}^{-2}_A(M_{n, 0}, \omega_A^\bullet)$
is at most the length of $\text{Ext}^{-2}_A(A/fA, \omega_A^\bullet)$
plus the length of the cokernel of
$$
t_0^n : 
\text{Ext}^{-2}_A(A/fA, \omega_A^\bullet)
\longrightarrow
\text{Ext}^{-2}_A(M_{n, 0}, \omega_A^\bullet)
$$
Via the isomorphism $M_0 \cong A/fA$ the map $t_0$ becomes
$g : A/fA \to A/fA$ for some $g \in A/fA$ and via the corresponding
isomorphisms $M_{n, 0} \cong A/fA$ the map $t_0^n$ becomes
$g^n : A/fA \to A/fA$. Thus the length of the cokernel above
is the length of the quotient of
$\text{Ext}^{-2}_A(A/fA, \omega_A^\bullet)$ by $g^n$.
Since $\text{Ext}^{-2}_A(A/fA, \omega_A^\bullet)$ is a finite $A$-module
with support $T$ of dimension $1$ and since $V(g) \cap T$
consists of the closed point by our choice of $t$
this length grows linearly in $n$ by
Algebra, Lemma \ref{algebra-lemma-support-dimension-d}.

\medskip\noindent
To finish the proof we prove the final assertion. Assume
$f \in \mathfrak m \subset A$ satisfies
(1), (2), (3), $A$ is $(S_2)$, and $\dim(A) \geq 4$.
Condition (1) implies $A$ is catenary, see
Dualizing Complexes, Lemma \ref{dualizing-lemma-universally-catenary}.
Then $\Spec(A)$ is equidimensional by Local Cohomology, Lemma
\ref{local-cohomology-lemma-catenary-S2-equidimensional}.
Thus $\dim(A_\mathfrak p) + \dim(A/\mathfrak p) \geq 4$
for every prime $\mathfrak p$ of $A$. Then
$\text{depth}(A_\mathfrak p) \geq \min(2, \dim(A_\mathfrak p))
\geq \min(2, 4 - \dim(A/\mathfrak p))$ and hence (4) holds.
\end{proof}

\begin{thebibliography}{9}
\bibitem{Kollar-map-pic}
Koll\'{a}r, J\'{a}nos.
\newblock \emph{Maps between local {P}icard groups}.
\newblock Algebr. Geom. 3 (4), 461--495 (2016).
\end{thebibliography}
\end{document}
